\documentclass[11pt]{article}

\usepackage[a4paper,margin=1in]{geometry}
\usepackage{amsmath,amssymb,amsthm}
\usepackage{booktabs}
\usepackage{multirow}
\usepackage{array}
\usepackage{graphicx}
\usepackage{hyperref}
\usepackage{xcolor}
\usepackage{microtype}
\usepackage{fancyvrb}
\usepackage{enumitem}
\usepackage[numbers,square]{natbib}
\usepackage{caption}
\usepackage{tabularx}

\hypersetup{
  colorlinks=true,
  linkcolor=blue!50!black,
  citecolor=blue!50!black,
  urlcolor=blue!50!black,
}

\newcommand{\hurts}{\text{$\dagger$}}
\newcommand{\helps}{\text{$\star$}}

\title{When the Map Hurts the Territory:\\
Articulated Internal-Model Prompts Degrade\\
Open-Weight LLM Sequence Prediction Across Scales}

\author{Mohamed A. Mabrok\\
\href{mailto:m.a.mabrok@gmail.com}{m.a.mabrok@gmail.com}}

\date{April 2026}

\begin{document}

\maketitle

\begin{abstract}
The Internal Model Principle (IMP) of \citet{francis1976internal} states
that any controller that perfectly tracks a class of exogenous signals
in steady state must contain, as a subsystem, a model of the signal
generator.  A natural translation of IMP to large language model (LLM)
agents is the prompting hypothesis: \emph{instructing an LLM to
articulate its internal model of the environment before acting should
improve closed-loop performance, especially as scale grows}.  We test
this hypothesis empirically on a $30$-task sequence-prediction
benchmark across three difficulty bands (easy, medium, hard) and four
open-weight model scales (Qwen2.5-7B, Qwen2.5-14B, Qwen2.5-32B,
Llama3.1-70B-q4) spanning roughly an order of magnitude in
parameter count, under three matched conditions: (i) a reactive
baseline that emits only the next integer, (ii) an IMP-style prompt
that asks the model to state a one-line hypothesis about $f(k)$
before predicting, and (iii) a zero-shot chain-of-thought variant.
Greedy decoding throughout.  We find: (a) articulating an internal
model \emph{never} significantly helps, on any difficulty band, at
any of the four scales; (b) on medium-difficulty tasks the
hypothesis-prompt significantly hurts at every tested scale
(paired-$t$ on $\log$-error, $p \in \{0.001, 0.010, 0.043,
0.025\}$; pooled Stouffer $p < 10^{-6}$); (c) the easy-band harm
decays \emph{monotonically} with scale (Spearman $\rho = -1.00$,
Pearson $r = -0.998$, $p = 0.002$ over the four scales) but does
not eliminate the medium-band harm, which persists at every scale;
and (d) catastrophic single-cell failures (e.g.\ Fibonacci with
reactive error~$0.00$ vs.\ hypothesis-prompt error~$71.17$ at $7$B,
geometric-2 with reactive~$0.00$ vs.\ hypothesis-prompt~$159.50$ at
$32$B) reveal a clear failure mode:
the model emits a wrong stated rule and then deterministically
computes forward from it.  These results are inconsistent with the
naive IMP-as-prompting reading and consistent with the meta-analytic
finding of \citet{sprague2024cot} that chain-of-thought helps mainly
on math/symbolic reasoning.  We discuss implications for treating
LLMs as controllers, list the experimental scope (one paradigm,
greedy decoding, integer targets, single seed), and pre-register
extensions.

\bigskip
\noindent\textbf{Keywords:} internal model principle, large language
models, sequence prediction, prompt engineering, chain-of-thought,
control theory, scaling.
\end{abstract}

\section{Introduction}
\label{sec:intro}

A theme of recent work in agentic AI is that LLMs increasingly play
the role of \emph{controllers} in closed loops: planners, tool-using
agents, multi-agent committees, and online predictors.  When a
controller is asked to track an exogenous reference, control theory
provides one of the few non-trivial sufficient/necessary conditions in
the field --- the Internal Model Principle, due to
\citet{francis1976internal}.  The IMP says, informally: \emph{to
achieve perfect asymptotic tracking of a class of disturbance or
reference signals, the controller must replicate, as an internal
subsystem, the dynamics that generate those signals}.

The IMP is a structural result about feedback architectures.  It is
silent about \emph{how} the model is encoded.  In a classical setting
the model lives in a state-space realization of the controller.  In
a learned setting it could live in network weights and activations.
In an LLM-agent setting it is tempting to read IMP as a prescription
about prompts: \emph{tell the LLM to write down its model of the
generator before acting, and tracking will improve}.  This reading is
attractive because it is operational --- one can implement
``maintain an internal model'' simply by editing the system prompt.
It is also the reading that practitioners have implicitly been
exploring for years under the banner of chain-of-thought
\citep{wei2022chain,kojima2022zero} and structured planning prompts.

This paper tests that reading directly.  We construct a clean
sequence-prediction game in which a hidden generator $f(k)$ produces
an integer target at each step $k$, the LLM observes the past and
emits a prediction $\hat y_k$, then the true $y_k$ is appended to the
history.  We compare three system prompts: a \emph{reactive} baseline
that asks for one integer with no scaffolding, an \emph{IM-Hypothesis}
prompt that explicitly invokes the IMP and instructs the model to
state a one-line hypothesis about $f(k)$ before predicting, and an
\emph{IM-CoT} zero-shot chain-of-thought prompt.  We run the same
$30$-task benchmark, balanced across three difficulty bands, on three
open-weight models spanning roughly an order of magnitude in scale.

The naive IMP-as-prompting prediction is straightforward:
articulating an internal model should help --- not necessarily
uniformly, but at least in the regime where the underlying pattern
is structured enough that having a model is informative.  This
prediction can be sharpened into three pre-registered hypotheses:
(H1) IMP prompts hurt small models on tasks they can already solve
reactively, but the harm decays with scale; (H2) IMP prompts help
small models on hard tasks (replicating the chain-of-thought
literature); (H3) the joint behaviour shows a significant
\emph{difficulty $\times$ condition $\times$ scale} interaction.

Our empirical findings refine H1, refute H2, and confirm H3 in a
modified form:
\begin{enumerate}[leftmargin=2em,itemsep=0pt,topsep=2pt]
\item Articulating an internal model \emph{never} significantly helps,
  on any of the twelve \emph{difficulty $\times$ scale} cells.
\item On \emph{medium}-difficulty tasks the IM-Hypothesis prompt
  significantly hurts at all four scales ($p$ values $0.001$, $0.010$,
  $0.043$, $0.025$; pooled Stouffer $Z = 5.06$, $p < 10^{-6}$).  The
  harm survives an order of magnitude of scaling.
\item The easy-band harm decays \emph{monotonically} with scale.
  The mean log-error gap (IM-Hypothesis $-$ reactive) shrinks from
  $+1.41$ at $7$B to $+1.08$ at $14$B to $+0.51$ at $32$B to $0.00$
  at $70$B-q4.  Spearman $\rho = -1.00$ ($p = 0$ at sample size $4$),
  Pearson $r = -0.998$ ($p = 0.002$).  This is a textbook clean
  scale-dependent decay --- but it does not transfer to the medium
  band, where the harm persists at every scale.
\item The harm does not so much decay across the board as
  \emph{migrate}: at $7$B and $14$B the IM-Hypothesis prompt damages
  \emph{easy} tasks most (because reactive performance there is
  excellent and there is room to fall); at $32$B and $70$B reactive
  performance on easy approaches the floor, so the same prompt
  instead damages \emph{medium} tasks where there is still slack.
\item Several individual cells exhibit a clean catastrophic failure:
  the reactive condition scores zero error while the IM-Hypothesis
  or IM-CoT condition produces large persistent error
  (e.g.\ Fibonacci on Qwen2.5-7B: $0.00$ vs.\ $71.17$;
  geometric-2 on Qwen2.5-7B and Qwen2.5-32B: $0.00$ vs.\ $\sim 160$;
  geometric-2 on Qwen2.5-14B with IM-CoT: $0.00$ vs.\ $384$;
  squares-minus-$k$ on Qwen2.5-32B with IM-CoT: $0.00$ vs.\ $68.67$).
  Inspection of the raw replies reveals a consistent mechanism ---
  a wrong rule is committed in the hypothesis or reasoning line, and
  integer predictions are then deterministically computed forward
  from the wrong rule.
\end{enumerate}

These observations are not a refutation of the IMP itself, which is a
mathematical theorem about feedback systems with explicit state
\citep{francis1976internal}.  They are evidence that the
\emph{prompting reading} of IMP --- ``ask the LLM to externalize an
internal model and use it'' --- is not a free lunch and can
actively backfire.  This is consistent with the broader picture
emerging from the chain-of-thought literature: the recent
meta-analysis of \citet{sprague2024cot} finds that CoT helps
predominantly on math/symbolic-execution tasks, with much smaller
gains elsewhere.  Our sequence-prediction setting is symbolic in
nature, which is precisely the regime where one might expect
articulation to help; that articulation \emph{hurts} here is a
non-trivial negative result.

The remainder of this paper is organized as follows.
Section~\ref{sec:imp} restates the classical IMP, distinguishing
mathematical content from prompting analogue.
Section~\ref{sec:related} positions the work within the
chain-of-thought, multi-agent failure, and LLM-as-controller
literatures.  Section~\ref{sec:method} gives the full benchmark
description --- exact prompts, exact tasks, exact metric, exact
decoding --- so the experiments are reproducible from the
description alone.  Section~\ref{sec:results} reports the main
table and the formal tests.  Section~\ref{sec:discussion} interprets
the failure mechanism.  Section~\ref{sec:limitations} states the
scope honestly: one task family, one decoding regime, integer
targets, single seed, three open models.  Section~\ref{sec:extensions}
lists pre-registered follow-up experiments that we believe are
necessary before any general claim can be made.

\section{The Internal Model Principle}
\label{sec:imp}

We restate the IMP in its classical form to keep the prompting
analogue under tight scrutiny.  Consider an LTI plant
\(P\) and an LTI controller \(K\) interconnected in a standard
unity-feedback loop, with reference $r$, output $y$, and tracking
error $e = r - y$.  Suppose the reference belongs to a class of
signals generated by an exosystem with state-space realization
\((A_e, B_e, C_e)\) --- for example, $A_e$ has poles at $\{0,
\pm j\omega_1, \pm j\omega_2, \dots\}$ for steps and persistent
sinusoids.  Then \citep[Theorem 1]{francis1976internal} establishes:
\emph{If the closed loop is internally stable and asymptotically
rejects the disturbance class generated by $(A_e, B_e, C_e)$, then
the dynamic modes of the exosystem must be reproduced inside the
loop} --- that is, the controller (or some part of the loop) must
contain a copy of the unstable modes of $A_e$.  This is a
\emph{necessity} result: stability plus tracking implies the
internal model.

Three observations bear on the prompting analogue.

\paragraph{The IMP is about modes, not language.}  The classical
result fixes a representation (state-space transfer, Laplace
domain) and the conclusion is about which poles appear in the loop.
Translating ``contains the model'' to ``writes down the model in
words'' is not implied by the theorem.  The implicit model of an
LLM, if any, lives in its weights and activations.  Articulating
that implicit model in tokens is a separate operation that need
not be representation-preserving and may be lossy.

\paragraph{The IMP is about exact asymptotic tracking, not
approximate prediction.}  The classical result is sharp under the
condition of \emph{zero} steady-state error.  An LLM agent doing
sequence prediction with finite horizon is not in the regime where
this sharp condition is operative; the relevant criterion is
average error, not exact tracking.

\paragraph{The IMP says ``contains'', not ``commits to''.}  An
exogenous signal class admitting a high-dimensional internal model
permits multiple realizations; the controller is free to maintain a
distribution over candidates, condition on data, and refine.  By
contrast, an LLM that has just emitted ``HYPOTHESIS:
$f(k) = ak + b$ with $a = 5, b = 1$'' has \emph{committed} to one
realization in the next-token sense: subsequent integer predictions
will be conditioned on that text, and the autoregressive decoding
will lock in the implied computation.

These three caveats already predict that the prompting analogue is
not entitled to the strength of the original theorem.  The empirical
question is whether, despite these caveats, the analogue gives a
useful heuristic.  Our answer is: not in the setup tested.

\section{Related Work}
\label{sec:related}

\paragraph{Chain-of-thought prompting.}
\citet{wei2022chain} demonstrate that prompting LLMs to produce
intermediate reasoning steps before answering improves accuracy on
arithmetic, commonsense, and symbolic-reasoning tasks --- but
chiefly on the largest models tested ($\geq$ 100B parameters in
their original setup).  \citet{kojima2022zero} simplify this to a
zero-shot ``Let's think step by step'' prompt, the exact phrasing
we use in our IM-CoT condition.  The recent meta-analysis of
\citet{sprague2024cot} surveyed over $100$ CoT papers and ran
controlled evaluations on $20$ datasets across $14$ models; their
headline finding is that CoT gives strong gains primarily on math
or logic, with much smaller gains on other tasks.  Our work fits
this picture: sequence prediction over deterministic integer
generators is a math/logic-style task, exactly the regime where
CoT should help; we observe instead that articulation hurts when
implemented as our paper's IM prompts.

\paragraph{Multi-agent LLM failures.}  Beyond single-agent
prompting, recent work investigates failure modes of multi-agent
LLM systems.  \citet{cemri2025why} introduce MAST, a taxonomy of
$14$ failure modes across $7$ frameworks built from rigorous
human annotation of $1{,}600+$ traces.  Their Category~I (system
design issues) includes failures driven by prompt structure and
role specification.  Our finding contributes a single-agent
specialization of that family: a system prompt that explicitly
asks for an articulated internal model can drive accuracy down on
a task family the same model would solve with no scaffolding.

\paragraph{Control-theoretic perspectives on LLMs.}
Several lines of recent work apply control-theoretic concepts to
neural systems --- stability, observability, feedback.  Our paper
is in this lineage but takes the opposite stance from much of it:
rather than borrowing a control-theoretic guarantee and asserting
that the LLM inherits it, we test, empirically and conservatively,
whether one of the most foundational such guarantees translates
\emph{at the prompting level}.  It does not.  This contributes a
data point for the broader question of which control-theoretic
results have meaningful prompt-level analogues and which require
working at the level of weights, activations, or training
objectives.

\section{Benchmark and Method}
\label{sec:method}

We describe the benchmark in full so the experiments are
reproducible from this section alone.  The implementation lives at
\url{https://github.com/Minds-R-Lab/phase_margin_llm}, file
\verb|pipeline/experiments/run_imp_benchmark.py|.

\subsection{Tasks}
\label{sec:tasks}

The benchmark contains $30$ deterministic integer-valued generators
$f \colon \mathbb{N} \to \mathbb{Z}$, balanced across three
difficulty bands of $10$ tasks each.  Difficulty was hand-assigned
according to the following criteria:

\textbf{Easy}: a no-reasoning baseline that simply echoes a recent
observation tends to be approximately correct (constants, slow
linear, low-period cycles).

\textbf{Medium}: identifying $f$ requires explicit pattern
recognition but, once identified, only one symbolic step is needed
to compute the next value (quadratic, longer periods, piecewise,
parity-modulated).

\textbf{Hard}: even with a correct hypothesis, the next-step
computation is non-trivial (recurrences, polynomial in $k$,
conditional rules, regime switches).

The full task list is given in Table~\ref{tab:tasks}.  All tasks use
$K = 14$ steps with the asymptotic error computed over the last
$6$, except \verb|h02_geometric_2| which is capped at $K = 12$ to
avoid integer overflow in the prompt.

\begin{table}[t]
\centering
\caption{The 30 benchmark tasks.  $f(k)$ is the deterministic
integer-valued generator; $K$ is the trajectory length; the last
$\textsc{asym}$ steps are used for the asymptotic error.  The
description column is intentionally short.  Tasks are addressed by
their two-letter prefix in the rest of the paper.}
\label{tab:tasks}
\small
\begin{tabularx}{\linewidth}{l l X r r}
\toprule
ID & Difficulty & Generator $f(k)$ & $K$ & \textsc{asym} \\
\midrule
\texttt{e01\_constant\_seven}      & easy   & $7$                                          & 14 & 6 \\
\texttt{e02\_constant\_zero}       & easy   & $0$                                          & 14 & 6 \\
\texttt{e03\_linear\_2k+3}         & easy   & $2k + 3$                                     & 14 & 6 \\
\texttt{e04\_linear\_5k}           & easy   & $5k$                                         & 14 & 6 \\
\texttt{e05\_linear\_desc\_3}      & easy   & $20 - 3k$                                    & 14 & 6 \\
\texttt{e06\_period2\_3\_8}        & easy   & cycle $[3,8]$                                & 14 & 6 \\
\texttt{e07\_period2\_10\_0}       & easy   & cycle $[10,0]$                               & 14 & 6 \\
\texttt{e08\_period3\_2\_6\_9}     & easy   & cycle $[2,6,9]$                              & 14 & 6 \\
\texttt{e09\_period4\_5\_10\_5\_0} & easy   & cycle $[5,10,5,0]$                           & 14 & 6 \\
\texttt{e10\_mod5}                 & easy   & $k \bmod 5$                                  & 14 & 6 \\
\midrule
\texttt{m01\_quadratic\_k2}        & medium & $k^2$                                        & 14 & 6 \\
\texttt{m02\_quadratic\_offset}    & medium & $(k - 1)^2$                                  & 14 & 6 \\
\texttt{m03\_period5\_pi}          & medium & cycle $[3,1,4,1,5]$                          & 14 & 6 \\
\texttt{m04\_period6\_e}           & medium & cycle $[2,7,1,8,2,8]$                        & 14 & 6 \\
\texttt{m05\_piecewise\_v}         & medium & $k$ if $k < 5$ else $10 - k$                 & 14 & 6 \\
\texttt{m06\_zigzag}               & medium & $k \cdot (-1)^k$                             & 14 & 6 \\
\texttt{m07\_step\_floor3}         & medium & $\lfloor k/3 \rfloor$                        & 14 & 6 \\
\texttt{m08\_triangle\_wave}       & medium & $|(k \bmod 6) - 3|$                          & 14 & 6 \\
\texttt{m09\_linear\_plus\_period} & medium & $k + 3 \cdot (k \bmod 2)$                     & 14 & 6 \\
\texttt{m10\_quadratic\_bounded}   & medium & $-k^2 + 10k$                                 & 14 & 6 \\
\midrule
\texttt{h01\_fibonacci}            & hard   & $f(0)=f(1)=1$, $f(k)=f(k{-}1)+f(k{-}2)$       & 14 & 6 \\
\texttt{h02\_geometric\_2}         & hard   & $2^k$ (capped to $4096$)                     & 12 & 4 \\
\texttt{h03\_squares\_minus\_k}    & hard   & $k(k - 1)$                                   & 14 & 6 \\
\texttt{h04\_triangle\_numbers}    & hard   & $k(k+1)/2$                                   & 14 & 6 \\
\texttt{h05\_pi\_digits}           & hard   & first 16 digits of $\pi$, cycled             & 14 & 6 \\
\texttt{h06\_e\_digits}            & hard   & first 16 digits of $e$, cycled               & 14 & 6 \\
\texttt{h07\_cubic\_mod100}        & hard   & $k^3 \bmod 100$                              & 14 & 6 \\
\texttt{h08\_conditional\_rule}    & hard   & $2k$ if $k \bmod 3 = 0$ else $k+1$           & 14 & 6 \\
\texttt{h09\_regime\_switch}       & hard   & cycle $[1,4]$ for $k<7$ then cycle $[5,2]$    & 14 & 6 \\
\texttt{h10\_sum\_two\_periods}    & hard   & $(k \bmod 3) + (k \bmod 5)$                  & 14 & 6 \\
\bottomrule
\end{tabularx}
\end{table}

\subsection{Conditions and Verbatim Prompts}
\label{sec:prompts}

Each task is run under three conditions, each defined by a single
system prompt.  We give the prompts verbatim because the central
finding of this paper depends on subtle wording differences.

\paragraph{Reactive (baseline).}
\begin{Verbatim}[fontsize=\small,frame=single]
You are a sequence-prediction agent.  At each step you will be
shown all observations so far in the form
`k=0: y=...; k=1: y=...; ...` and asked to predict the next value
y_k.  Reply with EXACTLY one integer on a line by itself.  No
commentary, no labels, no quotes.
\end{Verbatim}

\paragraph{IM-Hypothesis (the canonical IMP-as-prompt).}
\begin{Verbatim}[fontsize=\small,frame=single]
You are a sequence-prediction agent operating under the Internal
Model Principle.  At each step you will be shown all observations
so far and asked to predict the next value y_k.  You MUST first
state a one-line HYPOTHESIS about the function f(k) that generated
the observed values, then on a NEW LINE write your single integer
prediction for y_k.

Reply format (exactly two lines):
    HYPOTHESIS: <your one-line hypothesis>
    <integer prediction>
\end{Verbatim}

\paragraph{IM-CoT (zero-shot chain-of-thought variant).}
\begin{Verbatim}[fontsize=\small,frame=single]
You are a sequence-prediction agent.  At each step you will be
shown all observations so far and asked to predict the next value
y_k.  Let's think step by step: briefly reason about what rule
could be generating these values and why your candidate next
value is consistent with it.  After your reasoning, write your
single integer prediction for y_k as the LAST line of your reply,
by itself, with no labels or punctuation.  The integer on the
last line will be parsed automatically.
\end{Verbatim}

The user prompt is identical across all three conditions:
\begin{Verbatim}[fontsize=\small,frame=single]
Observations so far: k=0: y=...; k=1: y=...; ...

Predict y_k.
\end{Verbatim}

The reactive condition uses \verb|max_tokens = 96|, the IM-Hypothesis
condition uses \verb|max_tokens = 96|, the IM-CoT condition uses
\verb|max_tokens = 256| (because step-by-step replies are longer).
Decoding is greedy (\verb|temperature = 0|).

\subsection{Metric}
\label{sec:metric}

For each (task, condition, model) cell we run a single trajectory of
$K$ steps.  The agent observes the history, emits a reply, and the
last integer in the reply is parsed as the prediction $\hat y_k$.
The trajectory is then advanced by appending the \emph{true} $y_k$
(not the prediction), so the agent sees the same history regardless
of its own past errors --- the loop is open in the prediction
direction.  The asymptotic error of the cell is the mean absolute
error over the last $\textsc{asym}$ steps:
\[
  E_{\text{asym}} \;=\; \frac{1}{\textsc{asym}}
  \sum_{k = K - \textsc{asym}}^{K - 1}
  \big| \hat y_k - y_k \big|.
\]
For statistical comparisons we work with $\log(E_{\text{asym}} +
\varepsilon)$ with $\varepsilon = 0.5$ to keep the metric defined
when one or both conditions score zero.

For each pair of conditions $(a, b)$ and each difficulty band, we
form the ten differences
$\log(E_a + \varepsilon) - \log(E_b + \varepsilon)$ over the
$10$ tasks in that band and run a paired $t$-test.  A negative
$t$-statistic means condition $b$ produces lower log-error on
average; we report $b$ as ``hurts vs.\ $a$'' when this is
significant at $\alpha = 0.05$.

\subsection{Models and Decoding}
\label{sec:models}

We test four open-weight model scales served via Ollama on a
single H100 GPU:
\begin{itemize}[leftmargin=2em,itemsep=0pt,topsep=2pt]
\item \textbf{Qwen2.5-7B-Instruct} (\texttt{qwen2.5:7b}) -- $7.6$B
  parameters, $\sim 190$ s wall-time for the full benchmark.
\item \textbf{Qwen2.5-14B-Instruct} (\texttt{qwen2.5:14b}) --
  $14.7$B parameters, $\sim 580$ s.
\item \textbf{Qwen2.5-32B-Instruct} (\texttt{qwen2.5:32b}) --
  $32.5$B parameters, $\sim 1140$ s.
\item \textbf{Llama3.1-70B-Instruct, q4\_K\_M}
  (\texttt{llama3.1:70b-instruct-q4\_K\_M}) -- $70$B parameters,
  $4$-bit quantized, $\sim 9300$ s.
\end{itemize}
The four scales span roughly an order of magnitude in parameter
count.  All four are open-weight; the 70B is the only quantized
model in the set.  Frontier-scale closed models (Claude, GPT, Gemini)
are not tested in this paper; see Section~\ref{sec:extensions}.

Decoding is greedy in all cells (\verb|temperature = 0|, single
seed); this controls within-cell variance to zero, but at the cost
of within-task variance estimation.  Section~\ref{sec:limitations}
discusses this design choice and its consequences.

\section{Results}
\label{sec:results}

We report (i) the main mean-error table by difficulty $\times$
condition $\times$ scale, (ii) paired $t$-tests at $\alpha = 0.05$,
and (iii) the most-striking single-cell catastrophes that anchor
the failure-mode interpretation.

\subsection{Main Table}

Table~\ref{tab:main} reports the mean asymptotic error over the
ten tasks in each difficulty band, broken out by model scale and
prompting condition.  Numbers are mean
absolute error $|\hat y - y|$ on integer-valued targets, averaged
over the last six steps of a 14-step trajectory and then over the
ten tasks in the difficulty band.

\begin{table}[t]
\centering
\caption{Mean asymptotic error per (model, difficulty, condition).
Lower is better.  Bold marks cells where a paired $t$-test on
log-error against the same row's reactive baseline is significant
at $p < 0.05$ in the \emph{hurts} direction;
$\hurts$ marks the same significance.  No cell shows the help
direction at $p < 0.05$.}
\label{tab:main}
\small
\begin{tabular}{ll rrr}
\toprule
Model & Difficulty & Reactive & IM-Hypothesis & IM-CoT \\
\midrule
\multirow{3}{*}{Qwen2.5-7B}
  & easy   & $0.15$ & $\mathbf{6.08}^{\hurts}$  & $0.57$ \\
  & medium & $1.58$ & $\mathbf{5.22}^{\hurts}$  & $1.90$ \\
  & hard   & $5.58$ & $30.03$                   & $7.05$ \\
\addlinespace
\multirow{3}{*}{Qwen2.5-14B}
  & easy   & $0.00$ & $\mathbf{1.72}^{\hurts}$  & $\mathbf{1.32}^{\hurts}$ \\
  & medium & $0.43$ & $\mathbf{9.22}^{\hurts}$  & $\mathbf{2.57}^{\hurts}$ \\
  & hard   & $5.40$ & $5.43$                    & $45.12$ \\
\addlinespace
\multirow{3}{*}{Qwen2.5-32B}
  & easy   & $0.00$ & $0.75$                    & $0.25$ \\
  & medium & $1.73$ & $\mathbf{2.48}^{\hurts}$  & $2.30$ \\
  & hard   & $5.57$ & $20.85$                   & $\mathbf{13.53}^{\hurts}$ \\
\addlinespace
\multirow{3}{*}{Llama3.1-70B-q4}
  & easy   & $0.00$ & $0.00$\,\textsuperscript{(tied)}  & $0.02$ \\
  & medium & $0.37$ & $\mathbf{0.97}^{\hurts}$  & $\mathbf{1.83}^{\hurts}$ \\
  & hard   & $6.77$ & $8.03$                    & $11.25$ \\
\bottomrule
\end{tabular}
\end{table}

The shape of the table is the central empirical content of this
paper.  Reading column-wise: (a) reactive is the strongest column on
all three difficulty bands at $14$B and $70$B, and on easy and
medium at $7$B; (b) IM-Hypothesis is the worst column on $7$B in
all three bands and on $14$B medium; (c) IM-CoT is the worst column
on $14$B hard and on $70$B easy and medium.  Reading row-wise: (d)
both IM conditions inflate the mean across difficulty bands relative
to reactive, often by an order of magnitude or more.  Reading by
scale: (e) the easy column moves toward zero as scale grows for all
three conditions, but reactive reaches the floor first.

\subsection{Statistical Tests}

Table~\ref{tab:ttests} reports the paired-$t$ statistics on
$\log(E + 0.5)$ for each (scale, difficulty, IM condition) triple
against the reactive baseline.  $n = 10$ for every test.

\begin{table}[t]
\centering
\caption{Paired $t$-tests on log-error, reactive vs.\ each IM
condition, by scale and difficulty.  Each test has $n = 10$ tasks.
The $t$-statistics shown are the negation of the
log-difference test reported in the abstract: positive $t$ here
means the IM condition's log-error exceeds reactive's, i.e.\ IM
hurts.  At $\alpha = 0.05$, $11$ of $23$ testable cells are
significant in the hurt direction; zero in the help direction;
one cell (Llama-70B easy IM-Hyp) is not testable due to a
zero-variance tie at the floor.}
\label{tab:ttests}
\small
\begin{tabular}{l l rr rr}
\toprule
Model & Difficulty
  & \multicolumn{2}{c}{vs.\ IM-Hyp}
  & \multicolumn{2}{c}{vs.\ IM-CoT} \\
\cmidrule(lr){3-4}\cmidrule(lr){5-6}
 & & $t$ & $p$ & $t$ & $p$ \\
\midrule
\multirow{3}{*}{Qwen2.5-7B}
  & easy   & $+3.25$ & $\mathbf{0.010}^{\hurts}$ & $+1.38$ & $0.20$  \\
  & medium & $+4.73$ & $\mathbf{0.001}^{\hurts}$ & $+1.30$ & $0.23$  \\
  & hard   & $+2.20$ & $0.055$                   & $+1.53$ & $0.16$  \\
\multirow{3}{*}{Qwen2.5-14B}
  & easy   & $+3.38$ & $\mathbf{0.008}^{\hurts}$ & $+2.44$ & $\mathbf{0.037}^{\hurts}$ \\
  & medium & $+3.28$ & $\mathbf{0.010}^{\hurts}$ & $+4.58$ & $\mathbf{0.001}^{\hurts}$ \\
  & hard   & $+1.37$ & $0.20$                    & $+1.77$ & $0.11$  \\
\multirow{3}{*}{Qwen2.5-32B}
  & easy   & $+1.85$ & $0.097$                   & $+1.00$ & $0.34$  \\
  & medium & $+2.35$ & $\mathbf{0.043}^{\hurts}$ & $+1.13$ & $0.29$  \\
  & hard   & $+1.58$ & $0.15$                    & $+2.52$ & $\mathbf{0.033}^{\hurts}$ \\
\multirow{3}{*}{Llama3.1-70B-q4}
  & easy   & --      & --\,(tied at floor)        & $+1.00$ & $0.34$  \\
  & medium & $+2.68$ & $\mathbf{0.025}^{\hurts}$ & $+2.77$ & $\mathbf{0.022}^{\hurts}$ \\
  & hard   & $+0.60$ & $0.56$                    & $+1.50$ & $0.17$  \\
\bottomrule
\end{tabular}
\end{table}

\paragraph{Pooled significance on the medium band.}
The IM-Hypothesis condition hurts at every tested scale on the
medium-difficulty band.  Combining the four scale-specific
two-sided $p$-values via Stouffer's method (with sign of $t$ used
to orient the $z$-transform) yields $Z = 5.06$, $p < 10^{-6}$.
Stouffer requires independence across the combined tests; here
the four runs are independent in their stochasticity (different
models, separate cells), so the combined statement is sound.

\paragraph{Cross-scale trend on the easy band.}
The mean log-error gap of IM-Hypothesis vs.\ reactive on the easy
band is $+1.41, +1.08, +0.51, 0.00$ across Qwen-7B, 14B, 32B and
Llama-70B-q4 respectively.  Treating these four points as a
function of $\log(\text{parameter count})$ gives Spearman
$\rho = -1.00$ (perfect monotone decay; $p$ formally zero at
$n = 4$) and Pearson $r = -0.998$ ($p = 0.002$, two-sided).
The same trend test for the medium band gives Spearman $-0.80$
($p = 0.20$) and Pearson $-0.92$ ($p = 0.076$): a directionally
consistent decay that is not yet significant at four points but
falls outside the symmetry interval.

The IM-Hypothesis condition significantly hurts at all three scales
on medium-difficulty tasks ($p$ values $0.001$, $0.010$, $0.025$).
The IM-CoT condition does not hurt at $7$B (the smallest model is too
poorly-instructed to follow the CoT scaffolding tightly enough to
lock in a specific wrong rule), but does hurt on easy and medium at
$14$B and on medium at $70$B.  The hard band shows no significant
effect in either direction at any scale; reactive on hard is bad
enough that ratio-based tests have low power and the $t$-statistics
land in the $|t| < 2$ region.

The pattern that the harm \emph{migrates} rather than \emph{decays}
is most clearly visible by reading the IM-Hyp $p$-value column down
each scale block: $7$B has its smallest $p$-values on easy/medium;
$14$B keeps the hurt on easy/medium; $70$B loses the easy hurt only
because its reactive is at the floor.  In all three scales, the
\emph{medium} column is significantly hurt.

\subsection{Mechanism: Single-Cell Catastrophes}
\label{sec:mechanism}

The cleanest evidence about \emph{why} the IM prompt hurts comes
from individual cells in which the reactive condition achieves
zero asymptotic error and the IM-Hypothesis condition does not.
Such cells rule out the explanation ``the model just cannot do
the task.''  Table~\ref{tab:catas} lists the most-striking cases.

\begin{table}[t]
\centering
\caption{Catastrophic single-cell errors with reactive at floor.
The model demonstrably has the capability to predict the sequence
(reactive error $\leq 0.5$); inserting the IM scaffolding moves the
error into the tens or hundreds.  The phenomenon appears at every
tested scale --- including $32$B --- and on at least one
catastrophic cell (\texttt{h02\_geometric\_2}) under IM-Hyp it
reproduces almost numerically across the $7$B and $32$B scales.}
\label{tab:catas}
\small
\begin{tabular}{ll rrr}
\toprule
Model & Task & Reactive & IM-Hyp & IM-CoT \\
\midrule
Qwen2.5-7B  & \texttt{e04\_linear\_5k}        & $0.00$ & $41.67$  & $0.00$  \\
Qwen2.5-7B  & \texttt{h01\_fibonacci}         & $0.00$ & $71.17$  & $0.00$  \\
Qwen2.5-7B  & \texttt{h02\_geometric\_2}      & $0.00$ & $160.00$ & $0.00$  \\
Qwen2.5-7B  & \texttt{m02\_quadratic\_offset} & $0.00$ & $11.83$  & $3.83$  \\
Qwen2.5-14B & \texttt{m02\_quadratic\_offset} & $0.00$ & $66.17$  & $7.83$  \\
Qwen2.5-14B & \texttt{h02\_geometric\_2}      & $0.00$ & $0.00$   & $384.00$ \\
Qwen2.5-14B & \texttt{h03\_squares\_minus\_k} & $0.00$ & $5.33$   & $21.00$ \\
Qwen2.5-32B & \texttt{m02\_quadratic\_offset} & $0.00$ & $8.83$   & $16.50$ \\
Qwen2.5-32B & \texttt{m06\_zigzag}            & $0.00$ & $7.33$   & $0.00$  \\
Qwen2.5-32B & \texttt{h02\_geometric\_2}      & $0.00$ & $159.50$ & $0.00$  \\
Qwen2.5-32B & \texttt{h03\_squares\_minus\_k} & $0.00$ & $0.00$   & $68.67$ \\
Qwen2.5-32B & \texttt{h04\_triangle\_numbers} & $0.00$ & $1.17$   & $11.00$ \\
\bottomrule
\end{tabular}
\end{table}

Inspection of the raw replies in these cells reveals a consistent
mechanism: the model emits a wrong stated rule on the
\verb|HYPOTHESIS:| line --- a slope-by-one error, an off-by-one
quadratic, or a wrong recurrence --- and then on the next line
\emph{deterministically computes forward from the wrong rule},
producing predictions that drift further from the true
trajectory as $k$ grows.  The error is therefore not noise; it is
the hypothesis-text \emph{committing} the autoregressive decoder
to a particular forward computation.  When the same model is run
under the reactive prompt on the same task, no such commitment is
emitted and the implicit pattern-completion machinery (presumably
weights tuned for next-token prediction on natural-text sequences
that include integer patterns) suffices to predict correctly.

In control-theoretic terms: the IM-Hypothesis prompt does not
expose the implicit internal model that the LLM already has; it
forces the decoder to generate a \emph{separate}, articulated
internal model in tokens, which then becomes the operative model
for downstream prediction.  The articulated model is, in many
cells, worse than the implicit one.  This is the IM analogue of a
classical control failure: the controller and the disturbance
generator have the same modes, but the controller's
\emph{realization} of those modes is wrong, leading to negative
interference rather than tracking.

\section{Discussion}
\label{sec:discussion}

\subsection{Why does the implicit model outperform the articulated one?}

We offer one mechanism above; here we list three further candidate
explanations and the evidence for and against each.

\textbf{Hypothesis 1: tokenized commitment.}  The act of writing
\verb|HYPOTHESIS: <rule>| on line one creates a textual prefix on
which subsequent integer generation is conditioned.  If the rule
is wrong, the integer is forced toward the rule.  This matches
the catastrophic cells in Table~\ref{tab:catas}: the failures are
\emph{deterministic} given the wrong rule, not noisy.

\textbf{Hypothesis 2: distributional shift in the prompt.}  The
IM-Hypothesis prompt phrases the task as one in which the model
must do hypothesis testing about $f(k)$.  This may pull the
distribution of model behaviour toward training data resembling
mathematics homework or symbolic-regression problems --- contexts
where wrong intermediate symbolic claims are common.  The reactive
prompt, by contrast, puts the model in a vanilla-completion
regime that is much closer to the next-integer-in-a-list pattern
the model has seen in training data.  We have not yet probed this
directly; it is a concrete experiment for follow-up
(Section~\ref{sec:extensions}).

\textbf{Hypothesis 3: token budget pressure.}  IM-CoT gives more
\verb|max_tokens| than reactive, but its hurt is comparable to or
weaker than IM-Hypothesis.  This makes ``the IM condition has
\emph{too few} tokens to express the right rule'' an unlikely
sole explanation.  The IM-Hypothesis prompt also does not have a
token-budget problem: it consumes only one extra line.

\subsection{Comparison with the chain-of-thought literature}

\citet{wei2022chain} demonstrate that CoT helps on complex
multi-step reasoning, particularly for $\geq$ 100B parameter models.
Our largest tested model is $70$B-q4 (effectively smaller after
quantization), so we are formally below the regime where their
strongest gains were observed.  Nevertheless, our paradigm
(symbolic sequence prediction) is closer to the math/symbolic
tasks where CoT does help in the original setup.  Our finding ---
that CoT \emph{does not} help and the explicit-hypothesis variant
hurts --- is therefore not contradicted by the original CoT paper
but does narrow the regime of applicability further.  In the
language of \citet{sprague2024cot}, our task family
falls inside ``math/symbolic'' but the task structure (very short
integer sequences, deterministic generators, fixed
trajectory length) may not match the math problems on which CoT
is empirically strong.

\subsection{Implications for treating LLMs as controllers}

A take-away for the LLM-as-controller programme: \emph{the
existence of an implicit internal model in an LLM does not
license the prompt-level externalization of that model}.  In our
data, the implicit model embedded in network weights and
activations consistently outperforms the model the same network
articulates in tokens.  This has design implications for agentic
systems: prompting an LLM agent to ``maintain a world model'' or
to ``state its plan'' before acting cannot be assumed to
strengthen its control performance and may weaken it on
deterministic-tracking tasks.

\section{Limitations}
\label{sec:limitations}

We deliberately use this section to make the scope of our claim
precise.  All four limitations below are first-order and any one
of them could reverse the qualitative finding under different
choices.

\textbf{One paradigm, integer sequences.}  The benchmark is
single-agent sequence prediction with integer-valued, deterministic
generators.  Pattern types beyond what is in
Table~\ref{tab:tasks} are not tested.  Continuous-valued,
stochastic, or vector-valued tracking tasks may behave
differently.

\textbf{Greedy decoding, single seed.}  Setting
\texttt{temperature\,=\,0} eliminates within-cell variance but
suppresses any benefit
the IM-Hypothesis prompt might have under
\emph{self-consistency}-style decoding strategies, where multiple
sampled hypotheses are aggregated.  All single-seed claims should
be treated as upper bounds on the help an IMP-style prompt could
provide.

\textbf{Four open-weight scales, one quantization.}  We have
three Qwen2.5 sizes (7B, 14B, 32B) and one quantized Llama3.1
(70B, 4-bit).  No frontier-scale proprietary model is included,
and the largest model is 4-bit quantized so its effective capacity
is somewhere between $14$B and $70$B unquantized.  The
parameter-count span is roughly an order of magnitude, which is
enough to detect a monotone scale trend in the easy band but not
enough to extrapolate to two-orders-of-magnitude scaling.

\textbf{Three prompt phrasings.}  The harm we report is specific
to the exact prompts in Section~\ref{sec:prompts}.  Other
phrasings (e.g.\ ``maintain an internal model'' without the
explicit \verb|HYPOTHESIS:| line, or with explicit
self-correction instructions) might give different effects.

We have not run robustness sweeps over any of these dimensions.
For the level of generality the abstract claims, the data are not
yet sufficient.

\section{Extensions}
\label{sec:extensions}

We list, in priority order, the experiments needed to widen the
claim from ``hurts in this exact setup'' to ``hurts in general''
or to characterise the regime in which IMP-style prompting could
help.

\textbf{E1.\ Add a non-quantized 70B and additional intermediate
scales.}  Qwen2.5-32B has been added in this revision and confirms
the migration prediction: easy-band harm decays monotonically
(Spearman $\rho = -1.00$, Pearson $r = -0.998$, $p = 0.002$) while
medium-band harm persists.  The remaining gap is the unquantized
70B and ideally Qwen2.5-72B / Llama-3.1-405B for two-orders-of-
magnitude span.  Pre-registered prediction at larger scales: the
medium-band harm continues to shrink approximately log-linearly,
crossing into non-significance somewhere between $100$B and
$400$B parameters.

\textbf{E2.\ Self-consistency decoding.}  Re-run the IM-Hypothesis
condition at \verb|temperature = 0.7| with $N \in \{5, 10, 20\}$
sampled hypotheses per step, then either majority-vote the integer
or pick the integer matching the most-frequent hypothesis.  This
tests whether \emph{committing to one wrong rule} (our proposed
mechanism) is the actual driver: if averaging across hypotheses
recovers reactive performance, that is direct evidence for the
commitment mechanism.

\textbf{E3.\ Hypothesis quality probes.}  Programmatically classify
the generated \verb|HYPOTHESIS:| lines by correctness (does the
stated rule match the true generator?) and condition the analysis
on hypothesis correctness.  Predicted: cells where the hypothesis
is correct should match or beat reactive; cells where the
hypothesis is wrong should account for the entire IM-Hyp$-$reactive
gap.

\textbf{E4.\ Counterfactual prompts.}  Test IM-Hypothesis variants
with explicit invitation to revise the hypothesis when a prediction
fails (``If your previous prediction was wrong, update the
hypothesis''), and a self-consistency-aware variant that asks for
the top-$k$ hypotheses with weights.  These probe whether the
problem is the act of articulation or the absence of revision.

\textbf{E5.\ Continuous-valued targets.}  Replace integer
generators with continuous functions (sinusoids, slowly-varying
nonlinear maps) and ask the agent to emit floating-point
predictions.  IMP in classical control is most natural over
continuous signals; the prompting analogue may behave differently
when the targets do not invite token-by-token integer
commitment.

\textbf{E6.\ Multi-agent setting.}  In MAS, the IMP analogue is
that one agent plays the role of an explicit world-model maintainer
and the other plays the controller.  The MAST taxonomy of
\citet{cemri2025why} suggests this architecture frequently
fails for reasons distinct from single-agent prompting.  Running
our benchmark with a two-agent (modeller, predictor) split, where
the modeller's output is the predictor's prefix, would test
whether the failure mode we observe is purely a same-context
phenomenon or persists across agent boundaries.

\textbf{E7.\ Architecture probes.}  The mechanism we propose
(``tokenized commitment'') is testable via training-time
intervention: a model fine-tuned to emit hypotheses
\emph{conditional} on the integer prediction (i.e., generate the
integer first, then articulate the rule) should not exhibit the
forward-from-wrong-rule failure.  This requires a small SFT run
and is the most ambitious extension on this list.

\textbf{E8.\ Frontier-model replication.}  Run the same benchmark
on Claude Opus, GPT-4o, Gemini, and similar.  If the harm
disappears at very large scale, the headline becomes
``Below frontier scale, IMP-as-prompt hurts.''  If it persists,
the claim generalizes.  This experiment is straightforward but
incurs API costs and was out of scope for the open-weight
infrastructure used here.

\section{Why these results are useful}
\label{sec:usefulness}

\textbf{For practitioners building LLM agents.}  ``Maintain an
internal model'' prompts are common in agentic-system templates.
Our data suggest these prompts can degrade performance on
deterministic tracking, sometimes catastrophically.  The cheap
pre-deployment check is to run the benchmark in this paper (or a
domain-matched variant) under both reactive and articulated
conditions and pick the better one.  We expect this to flip
designs in non-trivial cases.

\textbf{For researchers borrowing from control theory.}  Our
result is a calibration data point for the broader programme of
mapping classical control concepts onto LLM agents.  IMP is one of
the strongest results in classical control; if its prompt-level
prompts, observability scaffolds) deserve cautious empirical
treatment before being treated as design principles.

\textbf{For the prompt-engineering literature.}  We add a clean,
reproducible cell to the negative-effects-of-CoT corpus
\citep{sprague2024cot} that targets a regime --- short symbolic
sequence prediction --- where one would naively expect CoT to
shine.  The benchmark is small enough to run in $\sim 10$ minutes
on a $7$B model, which makes it suitable as an inclusion criterion
for prompting papers that claim general benefits.

\textbf{For multi-agent failure analysis.}  The harm we observe
is the single-agent reduction of a class of failures the MAST
taxonomy \citep{cemri2025why} catalogues at the system level.  The
benchmark gives a controlled probe for the single-agent component
of those failures, which can then be composed back to study how
they propagate when an articulated wrong model is consumed by a
downstream agent.

\section{Conclusion}
\label{sec:conclusion}

We tested the prompting analogue of the Internal Model Principle
on a $30$-task sequence-prediction benchmark across four
open-weight model scales spanning roughly an order of magnitude
($7$B, $14$B, $32$B, $70$B-q4).  On no (difficulty, scale) cell
does articulating an internal model significantly help.  On
medium-difficulty tasks the explicit hypothesis prompt
significantly hurts at every tested scale (pooled Stouffer
$p < 10^{-6}$).  The easy-band harm decays cleanly and
monotonically with scale (Spearman $\rho = -1.00$,
Pearson $r = -0.998$, $p = 0.002$), but it does not generalize:
the medium-band harm persists at every scale, and catastrophic
single-cell failures of the same kind appear at every scale,
including $32$B.  Catastrophic single-cell failures expose the
mechanism: the articulated rule \emph{commits} the autoregressive
decoder to a wrong forward computation that the same model would
not have produced when prompted reactively.  These results are
inconsistent with the naive reading of IMP as a free-prompting
heuristic for LLM agents and consistent with the broader picture
that articulation does not uniformly improve LLM control behaviour.
The benchmark, prompts, and analysis code are released so that
each of the limitations and pre-registered extensions in
Sections~\ref{sec:limitations} and~\ref{sec:extensions} can be
addressed by independent replications.

\paragraph{Code and data.}
\url{https://github.com/Minds-R-Lab/phase_margin_llm}, file
\verb|pipeline/experiments/run_imp_benchmark.py|.

\bibliographystyle{plainnat}
\begin{thebibliography}{9}

\bibitem[Cemri et al.(2025)]{cemri2025why}
Mert Cemri, Melissa Z. Pan, Shuyi Yang, Lakshya A. Agrawal, Bhavya
Chopra, Rishabh Tiwari, Kurt Keutzer, Aditya Parameswaran, Dan Klein,
Kannan Ramchandran, Matei Zaharia, Joseph E. Gonzalez, and Ion Stoica.
\newblock Why Do Multi-Agent {LLM} Systems Fail?
\newblock arXiv preprint \href{https://arxiv.org/abs/2503.13657}{arXiv:2503.13657}, 2025.

\bibitem[Francis and Wonham(1976)]{francis1976internal}
B.~A. Francis and W.~M. Wonham.
\newblock The internal model principle of control theory.
\newblock \emph{Automatica}, 12(5):457--465, 1976.
\newblock \href{https://doi.org/10.1016/0005-1098(76)90006-6}{doi:10.1016/0005-1098(76)90006-6}.

\bibitem[Kojima et al.(2022)]{kojima2022zero}
Takeshi Kojima, Shixiang Shane Gu, Machel Reid, Yutaka Matsuo, and
Yusuke Iwasawa.
\newblock Large language models are zero-shot reasoners.
\newblock In \emph{Advances in Neural Information Processing Systems
(NeurIPS)}, 2022.
\newblock arXiv preprint \href{https://arxiv.org/abs/2205.11916}{arXiv:2205.11916}.

\bibitem[Sprague et al.(2024)]{sprague2024cot}
Zayne Sprague, Fangcong Yin, Juan Diego Rodriguez, Dongwei Jiang, Manya
Wadhwa, Prasann Singhal, Xinyu Zhao, Xi Ye, Kyle Mahowald, and Greg
Durrett.
\newblock To {CoT} or not to {CoT}? {C}hain-of-thought helps mainly on
math and symbolic reasoning.
\newblock arXiv preprint \href{https://arxiv.org/abs/2409.12183}{arXiv:2409.12183}, 2024.
\newblock Published at ICLR 2025.

\bibitem[Wei et al.(2022)]{wei2022chain}
Jason Wei, Xuezhi Wang, Dale Schuurmans, Maarten Bosma, Brian Ichter,
Fei Xia, Ed Chi, Quoc Le, and Denny Zhou.
\newblock Chain-of-thought prompting elicits reasoning in large
language models.
\newblock In \emph{Advances in Neural Information Processing Systems
(NeurIPS)}, 2022.
\newblock arXiv preprint \href{https://arxiv.org/abs/2201.11903}{arXiv:2201.11903}.

\end{thebibliography}

\end{document}
